\documentclass[11pt]{article}

\usepackage[margin=1in]{geometry}
\usepackage{amsmath,amssymb,amsthm}
\usepackage{booktabs}
\usepackage{graphicx}
\usepackage{hyperref}
\usepackage{xcolor}
\usepackage{array}
\usepackage{authblk}
\usepackage{siunitx}
\usepackage{natbib}

\hypersetup{
  colorlinks=true,
  linkcolor=blue!60!black,
  citecolor=blue!60!black,
  urlcolor=blue!60!black
}

\title{Calibrated Bayesian Overlap Detection for Long-Read\\
Genome Assembly}

\author{Raja Giddi}
\affil{Independent}
\date{April 2026}

\begin{document}
\maketitle

\begin{abstract}
Modern long-read genome assemblers report overlap calls without calibrated
uncertainty: a score above a fixed threshold is treated as a genuine overlap,
and downstream layout and consensus algorithms lose all information about
confidence. We present BAWM, a Bayesian overlap detector that outputs posterior
probabilities via likelihood-ratio Bayes factors over a small set of
interpretable features, combined with a bivariate $t$-copula correction for
feature dependence. A frozen encoder, contrastively pretrained on a
noisy-overlap proxy, supplies a single semantic feature; two geometric
features (k-mer containment and minimiser chain coverage) supply the rest.
On the \textit{Escherichia coli} K-12 MG1655 reference ($4.6$\,Mb, simulated
PacBio HiFi, 30\,$\times$ coverage, 9{,}283 reads, 462{,}875 candidate pairs),
BAWM attains precision $0.847$, recall $0.751$, and expected calibration
error $\mathrm{ECE} = 0.004$ at an interpretable decision threshold. The
calibration is roughly an order of magnitude tighter than the same pipeline
on a 100\,kb synthetic benchmark ($\mathrm{ECE}=0.040$) and is preserved in
the presence of seven near-identical rRNA operons, where the posterior
correctly collapses to low confidence ($\overline{p}=0.011$) rather than
fabricating overconfident edges. Feature ablation reveals that three of six
originally proposed orthogonal features contribute no signal or inverted
signal and are safely dropped; the three remaining features carry the full
result.
\end{abstract}

\section{Introduction}

De novo genome assembly from long reads is the problem of reconstructing a
genome sequence $G$ from a set of overlapping noisy observations
$\{r_1, \ldots, r_N\}$. The dominant pipeline --- all-versus-all overlap
detection, string-graph construction, and consensus \citep{myers2005fragment,
li2016minimap, cheng2021hifiasm} --- treats overlap detection as
classification: a score (e.g., minimiser chain length, shared $k$-mer count,
or mapping quality) is thresholded to produce binary edges. Whatever
uncertainty exists in the score is discarded.

This discards information that downstream algorithms could use. A string
graph constructed from binary edges cannot distinguish a ``probably real''
edge (posterior $0.6$) from a ``certainly real'' edge (posterior $0.99$);
consensus algorithms cannot weight their evidence by overlap confidence;
structural-variant callers cannot propagate uncertainty into calls. Worse,
the scores most assemblers rely on are not calibrated probabilities --- a
score of $0.6$ does not mean the edge is genuine $60\%$ of the time.

We take the view that \emph{overlap detection is a Bayesian inference
problem}, and the natural output is a posterior probability
$p(y_{ij} = 1 \mid \text{features}(r_i, r_j))$ that is (i) calibrated, so
that a stated probability matches the empirical frequency, and (ii) built
from interpretable, auditable features rather than an end-to-end black box.
Calibration lets downstream algorithms treat the posterior as usable
evidence; interpretability lets practitioners diagnose why an overlap was or
was not called.

\paragraph{Contributions.}
\begin{enumerate}
    \item A Bayesian overlap-detection model combining a small set of
    interpretable features via likelihood-ratio Bayes factors, with an
    explicit bivariate $t$-copula correction for feature dependence.
    \item A contrastive-pretraining procedure with anti-collapse
    regularisation that yields a frozen encoder usable as a single
    well-behaved feature, without fine-tuning or a learned confidence head.
    \item Results on \textit{E.\ coli} K-12 MG1655: precision $0.847$,
    recall $0.751$, $\mathrm{ECE}=0.004$, with calibration preserved in rRNA
    operon regions that typically degrade overlap callers.
    \item A feature ablation showing that three of six proposed orthogonal
    features are uninformative or inverted, and that the model performance
    is carried by three features; the cleaner model also has tighter
    calibration.
\end{enumerate}

\section{Method}

\subsection{Overview}

Given a set of reads $\{r_1, \ldots, r_N\}$, we consider the
$\binom{N}{2}$ candidate pairs (in practice pre-filtered by a cheap
minimiser-sharing criterion to $O(N \bar d)$ pairs where $\bar d$ is the
mean degree; for \textit{E.\ coli} this yields $462{,}875$ candidates over
$9{,}283$ reads). For each candidate pair $(i, j)$ we extract a feature
vector $\mathbf{x}_{ij}$ and compute the log posterior odds under the
two-class model $y_{ij} \in \{0, 1\}$ (null / genuine overlap) via

\begin{equation}
    \log \frac{p(y_{ij} = 1 \mid \mathbf{x}_{ij})}{p(y_{ij} = 0 \mid \mathbf{x}_{ij})}
    = \underbrace{\log\frac{\pi}{1-\pi}}_{\text{prior}} +
      \underbrace{\sum_{k=1}^{K} w_k \log\frac{p(x_k \mid y=1)}{p(x_k \mid y=0)}}_{\text{feature Bayes factors}} +
      \underbrace{\log\frac{c_1(\mathbf{u})}{c_0(\mathbf{u})}}_{\text{copula correction}},
    \label{eq:logit}
\end{equation}

where $\pi$ is the prior probability of a genuine overlap among candidate
pairs (estimated from a calibration set), $w_k \ge 0$ are fixed feature
weights, $p(x_k \mid y)$ are class-conditional distributions over the $k$-th
feature, and $c_y(\mathbf{u})$ is a copula density accounting for residual
dependence between features that are not conditionally independent given
the class.

The posterior $p(y_{ij} = 1 \mid \mathbf{x}_{ij}) = \sigma(\text{logit})$ is
the primary output of the model.

\subsection{Features}

We extract three features per candidate pair:

\paragraph{Containment.}
Let $\mathcal{M}_i, \mathcal{M}_j$ be the $(w, k) = (5, 15)$ minimiser sets
of reads $i$ and $j$. Define
\begin{equation}
    c_{ij} = \frac{|\mathcal{M}_i \cap \mathcal{M}_j|}{\min(|\mathcal{M}_i|, |\mathcal{M}_j|)}.
\end{equation}
This is the classical shared-$k$-mer containment score.

\paragraph{Chain coverage.}
Shared minimisers alone do not imply a collinear overlap; a pair of reads
whose shared minimisers are scattered across the genome (a common repeat
signature) will have high $c_{ij}$ but no overlap. We compute the longest
collinear chain of shared minimisers via sparse dynamic programming with
gap penalty (standard seed-chain-extend), and report
\begin{equation}
    f_{\text{cov}}(i,j) = \frac{|\text{longest colinear chain}|}{\min(|\mathcal{M}_i|, |\mathcal{M}_j|)}.
\end{equation}
Chain coverage is high when shared minimisers are in the same order on both
reads (expected for genuine overlaps) and low when they are scattered
(expected for repeats).

\paragraph{Encoder cosine.}
A transformer-style encoder $E_\theta : \{A,C,G,T\}^* \to \mathbb{R}^{64}$
(detailed in \Cref{sec:encoder}) maps each read to a fixed-dimensional
latent $\mu_i = E_\theta(r_i)$. The semantic similarity feature is
\begin{equation}
    s_{ij} = \frac{\langle \mu_i, \mu_j \rangle}{\|\mu_i\|\,\|\mu_j\|}.
\end{equation}
$E_\theta$ is trained once (\Cref{sec:encoder}) and held frozen during
calibration and inference; no classification head is attached.

\paragraph{Remark: features we tried and dropped.}
We originally included three additional features --- a chain score ratio
(secondary-to-primary chain), an end-anchoring score
$\exp(-\text{overhang}/\tau)$, and a TF--IDF style ``unique containment''
using inverse $k$-mer document frequency. Feature ablation on the
calibration data (\Cref{sec:ablation}) revealed each to have zero or
inverted class separation at \textit{E.\ coli} scale. They are dropped from
the final model and discussed in \Cref{sec:ablation} and
\Cref{sec:discussion}.

\subsection{Class-conditional distributions}

For each feature $x_k$ we fit class-conditional densities by method of
moments on the calibration set:
\begin{itemize}
    \item \textbf{Containment}, \textbf{chain coverage}: Beta distributions
    $p(x_k \mid y) = \mathrm{Beta}(a_y^{(k)}, b_y^{(k)})$, appropriate for
    $[0,1]$-valued features.
    \item \textbf{Encoder cosine}: Gaussian $p(x_k \mid y) =
    \mathcal{N}(\mu_y^{(k)}, \sigma_y^{(k)})$, appropriate for cosine
    similarities of unnormalised latents. (The cosine distribution is in
    practice concentrated away from the $\pm 1$ boundary.)
\end{itemize}

Weights are set as $w_{\text{containment}} = w_{\text{chain cov}} = 1$
and $w_{\text{cosine}} = 0.5$; the lower weight on cosine reflects that it
partially overlaps with containment in the evidence it carries.

\subsection{Copula correction}
\label{sec:copula}

Containment and encoder cosine are derived from overlapping substrates:
containment counts shared $k$-mers, and the encoder is trained to put reads
that share $k$-mers near each other. The assumption of conditional
independence $p(\mathbf{x} \mid y) = \prod_k p(x_k \mid y)$ is therefore
violated, and na\"ively summing log-Bayes factors double-counts evidence.

We correct for this with a bivariate $t$-copula \citep{nelsen2006copula}.
Given the empirical CDFs $F_{c,y}, F_{s,y}$ of containment and cosine under
class $y$, the copula reparametrises
\begin{equation}
    u_c = F_c(x_c), \quad u_s = F_s(x_s), \quad (u_c, u_s) \in [0,1]^2,
\end{equation}
and models the joint density as
\begin{equation}
    p(x_c, x_s \mid y) = p(x_c \mid y) \, p(x_s \mid y) \cdot c_y(u_c, u_s),
\end{equation}
so that the log Bayes factor in \Cref{eq:logit} acquires a correction term
\begin{equation}
    \Delta_{\text{cop}} = \log c_1(u_c, u_s) - \log c_0(u_c, u_s).
\end{equation}
We use a $t$-copula with $\nu = 5$ degrees of freedom (heavy tails are
appropriate for correlated sequence features with occasional extreme
co-occurrence) and estimate the correlation parameter $\rho_y$ from the
Spearman correlation of $(x_c, x_s)$ in each class of the calibration set.
On \textit{E.\ coli} we obtain $\rho_1 = 0.065$ (genuine) and $\rho_0 =
0.004$ (null): the dependence is modest but real under the genuine
distribution, as expected.

\subsection{Encoder pretraining}
\label{sec:encoder}

The encoder is a small transformer (4 layers, hidden $d=128$, latent dim
$64$, $\approx 1.2$\,M parameters) with positional embeddings and a mean-pool
readout. It is trained with an InfoNCE-style contrastive objective using
noisy ground-truth overlaps as positives:
\begin{equation}
    \mathcal{L} = -\log \frac{\exp(\mu_a \cdot \mu_p / \tau)}{\exp(\mu_a \cdot \mu_p / \tau) + \sum_{n} \exp(\mu_a \cdot \mu_n / \tau)}.
\end{equation}
Three anti-collapse measures were necessary in preliminary experiments:
(i) orthogonality kick, an initial additive perturbation of the readout
layer; (ii) a spread regulariser
$\mathcal{L}_{\text{spread}} = \frac{1}{|B|^2}\sum_{i \ne j} |\hat\mu_i \cdot \hat\mu_j|$
on $\ell_2$-normalised latents within a batch; and (iii) a temperature
schedule $\tau: 0.5 \to 0.1$ over the first third of training. Without
these, the encoder collapses to a near-constant output and the cosine
feature loses all signal.

Training runs for $30$ epochs on an NVIDIA A10G GPU with batch size $16$
and $16$ negatives per anchor. The encoder is then frozen for calibration
and inference.

\subsection{Calibration}
\label{sec:calibration}

Calibration has three outputs:
\begin{enumerate}
    \item the prior $\pi = n_{\text{genuine}}/n_{\text{pairs}}$;
    \item per-feature class-conditional parameters
    $(a_y^{(k)}, b_y^{(k)})$ or $(\mu_y^{(k)}, \sigma_y^{(k)})$;
    \item per-class copula correlations $\rho_y$.
\end{enumerate}

For \textit{E.\ coli} we use the full candidate set as the calibration set
with ground-truth labels derived from the simulator's true-position
annotations (see \Cref{sec:data}); in production use, calibration would be
run on a held-out reference genome. No neural network is trained in the
calibration phase.

\subsection{Decision threshold}
\label{sec:threshold}

The default decision threshold at log-odds $0$ corresponds to
$\pi$-weighted cost-neutral prediction. For assembly, a false negative
(missing a genuine overlap) fragments the layout and tends to be more
costly than a false positive (a spurious edge), which downstream consensus
and graph-cleaning steps can filter. We therefore select the operating
threshold by sweep on the calibration set; equivalently, this is a shift
in the effective log-prior by $\Delta$, so that
\begin{equation}
    t^* = \arg\max_t F_1\left(t\right) \quad \text{subject to} \quad P(t) \ge 0.80,\;R(t) \ge 0.75.
\end{equation}
On \textit{E.\ coli} the optimum is at $t^* = -1.5$, an effective cost
ratio of $e^{1.5} \approx 4.5 : 1$ for false negatives versus false
positives. The decision boundary is stable: every $t \in [-2.25, -1.50]$
satisfies the two-sided constraint, so the result is not knife-edge.

\section{Data}
\label{sec:data}

\paragraph{Reference.} \textit{Escherichia coli} K-12 MG1655 (GenBank
NC\_000913.3, $4{,}641{,}652$\,bp, circular). The genome contains seven
near-identical rRNA operons (rrnA--rrnH), each approximately $5.5$\,kb long,
which are the dominant source of repeat-induced overlap ambiguity in
bacterial assembly.

\paragraph{Reads.} PacBio HiFi reads simulated at $30\times$ coverage with
lognormal length distribution (mean $15$\,kb, sd $3$\,kb), a homopolymer-aware
error model ($\epsilon_{\text{sub}} = 0.003$, $\epsilon_{\text{indel}} =
0.001$, elevated at homopolymer runs), and uniform start positions. Yields
$9{,}283$ reads.

\paragraph{Ground-truth labels.} A pair $(i, j)$ is genuine if the true
genomic intervals of the two reads overlap by at least $5$\,kb (after
handling circular wraparound and strand). A spatial bin index over read
start positions reduces the $\binom{N}{2}$ label computation to linear in
$N$ without approximation. Total: $462{,}875$ candidate pairs, $6{,}638$
genuine ($\pi = 0.0143$).

\paragraph{Repeat-region subset.} For the rRNA operon analysis we define the
``repeat subset'' as all candidate pairs where at least one read's true
start position falls inside one of the seven rRNA operon intervals. This
gives $8{,}629$ pairs.

\section{Results}

\subsection{Main result}

\begin{table}[h]
\centering
\caption{Overall performance on \textit{E.\ coli} K-12 MG1655. BAWM-6f is
the model with all six originally proposed features; BAWM-3f drops the
three uninformative/inverted features (see \Cref{sec:ablation}). $t^*$ is
the log-odds threshold at which $F_1$ is maximised subject to
$P \ge 0.80, R \ge 0.75$. ECE is expected calibration error over 10 equal-width bins.}
\label{tab:main}
\begin{tabular}{lccccc}
\toprule
Model & Threshold & Precision & Recall & $F_1$ & ECE \\
\midrule
BAWM-6f ($t = 0$) & $0$    & $0.796$ & $0.751$ & $0.773$ & $0.004$ \\
BAWM-3f ($t = 0$) & $0$    & $0.905$ & $0.695$ & $0.786$ & $0.004$ \\
\textbf{BAWM-3f ($t^*$)} & $-1.5$ & $\mathbf{0.847}$ & $\mathbf{0.751}$ & $\mathbf{0.796}$ & $\mathbf{0.004}$ \\
\bottomrule
\end{tabular}
\end{table}

BAWM-3f at $t^* = -1.5$ reaches precision $0.847$, recall $0.751$, $F_1 =
0.796$, and expected calibration error $0.004$ on $462{,}875$ candidate
pairs. All three metrics clear the targets set \textit{a priori}
($P > 0.80$, $R > 0.75$, $\mathrm{ECE} < 0.12$) with margin.

\subsection{Calibration}
\label{sec:calibration-results}

\begin{figure}[h]
\centering
\includegraphics[width=0.95\linewidth]{../results/calibration_3feat.png}
\caption{Left: reliability diagram. Predicted confidence (x-axis) vs.\
empirical fraction of genuine overlaps (y-axis), binned into 10 bins.
Dashed line is perfect calibration. Right: confidence distribution split by
true label; note the clean separation and the near-absence of null pairs at
high confidence.}
\label{fig:calibration}
\end{figure}

The reliability diagram (\Cref{fig:calibration}, left) tracks the perfect-calibration line within bin-sampling noise across the full range
$[0, 1]$. The confidence distribution (\Cref{fig:calibration}, right) is
strongly bimodal: null pairs accumulate at very low confidence (mean
log-odds $-10.0$, sd $1.95$), genuine pairs at high confidence (mean
$+7.2$, sd $10.8$), with a sparse overlap region that sets the
precision-recall trade-off.

For context, the same pipeline on a 100\,kb synthetic benchmark (Paper 1,
same features minus chain coverage, single-feature BF without copula)
yielded ECE $= 0.040$. The \textit{E.\ coli} result is an order of
magnitude tighter despite the genome being $46\times$ larger and the
candidate set $\sim 10\times$ larger, driven primarily by the copula
correction (\Cref{sec:copula-effect}) rather than by feature engineering.

\subsection{Repeat-region behaviour}

rRNA operons are the canonical failure mode for bacterial overlap
detection. \Cref{tab:repeats} reports performance on the repeat subset.

\begin{table}[h]
\centering
\caption{Performance on the repeat subset (pairs with at least one read in
an rRNA operon). Mean confidence is the average posterior
$p(y=1\mid\mathbf{x})$ over the subset.}
\label{tab:repeats}
\begin{tabular}{lcccc}
\toprule
Subset & $n$ pairs & Precision & Recall & Mean conf. \\
\midrule
All pairs        & $462{,}875$ & $0.847$ & $0.751$ & $0.014$ \\
rRNA operon subset & $\phantom{00}8{,}629$ & $0.700$ & $0.656$ & $0.011$ \\
\bottomrule
\end{tabular}
\end{table}

Mean posterior confidence on rRNA pairs is $0.011$, below the $0.014$ base
rate: the posterior is \emph{correctly under-confident} in repeat regions,
which is the behaviour any downstream graph-cleaning or layout algorithm
wants. Precision drops to $0.70$ (vs.\ $0.85$ global), recall to $0.66$.
Crucially, calibration is preserved: the $0.70$ precision is what the
model predicts, not a silent failure.

\subsection{Threshold sensitivity}

\Cref{fig:threshold} plots precision, recall, and $F_1$ as functions of
the log-odds threshold. The region $t \in [-2.25, -1.50]$ satisfies the
two-sided constraint $P \ge 0.80, R \ge 0.75$, indicating a stable
operating region rather than a knife-edge optimum.

\subsection{Feature ablation}
\label{sec:ablation}

\begin{table}[h]
\centering
\caption{Per-feature class separation (mean under genuine minus mean under
null) on the \textit{E.\ coli} calibration set. Positive separation =
feature discriminates in the expected direction. The three features below
the rule contribute no signal or inverted signal and are dropped from the
final model.}
\label{tab:ablation}
\begin{tabular}{lrrr}
\toprule
Feature & Genuine mean & Null mean & Separation \\
\midrule
Containment           & $0.257$ & $0.133$ & $+0.124$ \\
Chain coverage        & $0.044$ & $0.014$ & $+0.030$ \\
Encoder cosine        & $0.784$ & $0.749$ & $+0.035$ \\
\midrule
Chain score ratio     & $0.757$ & $0.916$ & $-0.159$ \\
End-anchoring         & $0.001$ & $0.002$ & $-0.001$ \\
Unique containment    & $0.041$ & $0.051$ & $-0.011$ \\
\bottomrule
\end{tabular}
\end{table}

Three of the six originally proposed features are informative in the
direction assumed by the model (containment, chain coverage, encoder
cosine). Three are not:

\begin{itemize}
    \item \textbf{Chain score ratio} (secondary-to-primary chain ratio)
    has strongly \emph{inverted} separation. This is diagnosable: at
    \textit{E.\ coli} scale with $30\times$ coverage, null pairs that
    share minimisers in any repeat region produce long secondary chains,
    inflating the ratio; genuine pairs with clean single chains produce
    small ratios. A sign flip would rescue the feature, but the resulting
    signal is largely redundant with chain coverage.
    \item \textbf{End-anchoring} is dead: both classes have mean near
    $0.001$, because the decay scale of the $\exp(-\text{overhang}/\tau)$
    saturates at realistic overhangs (reads $\sim 15$\,kb, overlaps
    $\sim 5$\,kb, overhangs $\sim 10$\,kb). This is a design bug in the
    feature, not a fundamental issue with dovetail geometry.
    \item \textbf{Unique containment} (TF--IDF-weighted containment) has
    mildly inverted separation. We suspect the global document-frequency
    threshold is filtering too aggressively; a per-region tuning might
    rescue it, but we have not pursued this.
\end{itemize}

Dropping the three features has two effects: precision at $t = 0$ rises
from $0.796$ to $0.905$, and the null log-odds standard deviation drops
from $2.85$ to $1.95$. Recall at $t = 0$ drops from $0.751$ to $0.695$,
recovered by lowering the threshold to $t^* = -1.5$. ECE is essentially
unchanged ($0.004$ in both cases).

\subsection{Effect of the copula correction}
\label{sec:copula-effect}

\begin{table}[h]
\centering
\caption{Effect of the $t$-copula correction. ``Independent'' sets
$\rho_0 = \rho_1 = 0$; ``Corrected'' uses the estimated $\rho_1 = 0.065$,
$\rho_0 = 0.004$.}
\label{tab:copula}
\begin{tabular}{lccc}
\toprule
Model & Precision & Recall & ECE \\
\midrule
BAWM-3f, no copula correction & $0.836$ & $0.746$ & $0.012$ \\
BAWM-3f, $t$-copula corrected & $0.847$ & $0.751$ & $\mathbf{0.004}$ \\
\bottomrule
\end{tabular}
\end{table}

The copula correction is primarily a calibration fix, not a precision fix.
Removing it raises ECE by roughly $3\times$ while leaving precision and
recall approximately intact. This is the expected pattern: dependent
features bias the log-Bayes factor but typically in a smooth direction
that barely moves the decision boundary while damaging the probability
scale. On larger and more repeat-rich genomes this gap should widen.

\section{Discussion}
\label{sec:discussion}

\paragraph{What worked.}
The core recipe --- a small number of interpretable features combined by
likelihood-ratio Bayes factors, with copula correction for known feature
dependence --- generalises cleanly from 100\,kb to 4.6\,Mb. Calibration is
a property of the inference rule, not of the features, and survives the
scale jump; precision and recall are properties of the features, and the
\textit{E.\ coli} rRNA operons produce a visible but correctly quantified
degradation rather than a silent failure.

\paragraph{What did not work.}
Half of the originally proposed orthogonal features were uninformative or
inverted at \textit{E.\ coli} scale. The lesson is not that the features
were bad ideas, but that adding features without per-feature ablation is
dangerous: uninformative features are harmless in the infinite-data limit
but add variance to the log-Bayes factor at finite calibration-set size,
visibly widening the genuine-class distribution and degrading precision.
An honest methodology should ablate every feature.

\paragraph{Limitations.}
The ground truth used here is simulator-derived: true genomic coordinates
are known exactly, so calibration labels are noiseless. Real HiFi reads
(e.g., SRA SRR10971019) do not have this luxury; calibration on real data
would use a reference-guided alignment as an approximate gold standard,
introducing label noise. We have not yet tested whether the calibration
result survives this shift. We also have not compared against minimap2
\citep{li2016minimap} or similar production overlap callers on matched
data; their scores are not calibrated probabilities, so the comparison
would necessarily be on precision/recall at a fixed operating point and
not on ECE.

\paragraph{Toward a world model.}
The name BAWM (Bayesian Assembly as World Modeling) signals a longer-term
research programme in which the genome is treated as a latent variable and
the reads as noisy observations. The current paper is properly a calibrated
\emph{edge detector}, not a world model: it does not predict unobserved
reads, does not sample hypothetical genomes, and performs inference
pair-by-pair rather than globally over the assembly graph. The natural
next step is global inference via loopy belief propagation over the edge
graph, giving an edge marginal that depends on the whole graph rather than
two reads --- the first component of a proper world model.

\section{Related work}

\paragraph{Overlap detection.} Minimap2 \citep{li2016minimap} and the
overlap phase of hifiasm \citep{cheng2021hifiasm} are the canonical
long-read overlap callers, both based on minimiser chaining with hard score
thresholds. Neither outputs calibrated posterior probabilities.

\paragraph{Contrastive representation learning for biology.}
Contrastive objectives have been used to pretrain sequence encoders
(Nucleotide Transformer \citep{dalla2023nucleotide}, DNABert-2, HyenaDNA),
typically as the basis for supervised fine-tuning. Our use is less
ambitious: the encoder supplies a single feature in a larger Bayesian
model, never fine-tuned, and never the sole basis of a prediction.

\paragraph{Calibration in machine learning.}
Temperature scaling, Platt scaling, and isotonic regression
\citep{guo2017calibration} are the standard post-hoc calibration
techniques for deep classifiers. Our approach is structurally different:
calibration is achieved by construction through the Bayesian inference
rule and copula correction, not by a post-hoc transformation of
classifier logits.

\paragraph{Bayesian models in bioinformatics.}
Bayesian formulations have been applied to variant calling
\citep{li2011statistical}, phylogenetics, and read mapping, but to our
knowledge not to calibrated overlap detection for long-read assembly.

\section{Conclusion}

Long-read overlap detection can be formulated as a calibrated Bayesian
inference problem with a small number of interpretable features, and the
resulting posterior probabilities are usable by downstream algorithms in
ways that binary edges are not. On \textit{E.\ coli} K-12, the resulting
model reaches precision $0.847$, recall $0.751$, and
$\mathrm{ECE} = 0.004$, with calibration preserved in rRNA operon regions.
The calibration result survives a 46-fold scale jump from synthetic
benchmarks, is primarily driven by copula correction for feature
dependence rather than by feature engineering, and degrades gracefully
(with \emph{predicted} low confidence) in repeat regions rather than
failing silently. The natural next step is global inference over the edge
graph, turning the pair-wise posterior into a graph-wide one.

\section*{Code and data availability}
The implementation, simulator, and all artefacts used to produce the
\textit{E.\ coli} numbers in this paper are available at
\url{https://github.com/<anon>/bawm}.

\bibliographystyle{plainnat}
\begin{thebibliography}{99}

\bibitem[Myers(2005)]{myers2005fragment}
Eugene W.\ Myers.
\newblock The fragment assembly string graph.
\newblock \emph{Bioinformatics}, 21:\penalty0 ii79--ii85, 2005.

\bibitem[Li(2018)]{li2016minimap}
Heng Li.
\newblock Minimap2: pairwise alignment for nucleotide sequences.
\newblock \emph{Bioinformatics}, 34(18):\penalty0 3094--3100, 2018.

\bibitem[Cheng et al.(2021)]{cheng2021hifiasm}
Haoyu Cheng, Gregory T.\ Concepcion, Xiaowen Feng, Haowen Zhang, and Heng Li.
\newblock Haplotype-resolved de novo assembly using phased assembly graphs
  with hifiasm.
\newblock \emph{Nature Methods}, 18:\penalty0 170--175, 2021.

\bibitem[Nelsen(2006)]{nelsen2006copula}
Roger B.\ Nelsen.
\newblock \emph{An Introduction to Copulas}.
\newblock Springer, 2nd edition, 2006.

\bibitem[Dalla-Torre et al.(2023)]{dalla2023nucleotide}
Hugo Dalla-Torre et al.
\newblock The Nucleotide Transformer: Building and evaluating robust
  foundation models for human genomics.
\newblock \emph{bioRxiv}, 2023.

\bibitem[Guo et al.(2017)]{guo2017calibration}
Chuan Guo, Geoff Pleiss, Yu Sun, and Kilian Q.\ Weinberger.
\newblock On calibration of modern neural networks.
\newblock In \emph{ICML}, 2017.

\bibitem[Li(2011)]{li2011statistical}
Heng Li.
\newblock A statistical framework for SNP calling, mutation discovery,
  association mapping and population genetical parameter estimation from
  sequencing data.
\newblock \emph{Bioinformatics}, 27(21):\penalty0 2987--2993, 2011.

\end{thebibliography}

\end{document}
